\subsection{连通空间的积}

命题 8. 连通空间的每个积都是连通的。反之，若非空空间的一个积连通，则每个因子都连通。

设 \(\mathbf{X} = \prod_{i \in I} \mathbf{X}_i\) 是拓扑空间的一个积。若诸 \(\mathbf{X}_i\) 非空，则有 \(\mathbf{X}_i = \text{pr}_i\mathbf{X}\)（对每个 \(i \in I\)）；因而若 \(\mathbf{X}\) 连通，则诸 \(\mathbf{X}_i\) 也连通（第 2 目，命题 4）。反之，设诸 \(\mathbf{X}_i\) 都连通而 \(\mathbf{X}\) 不连通。由第 2 目命题 5，存在一个连续满射 \(f: \mathbf{X} \to \mathbf{X}'\)，其中 \(\mathbf{X}'\) 是含有多于一个点的离散空间。设 \(a = (a_i)\) 是 \(\mathbf{X}\) 的任意一点，\(x\) 是任意指标；部分映射 \(f_x: \mathbf{X}_x \to \mathbf{X}'\) 由 \(f_x(x) = f((y_i))\) 定义，其中 \(y_x = x\)，而 \(y_i = a_i\)（若 \(i \neq x\)），它在 \(\mathbf{X}_x\) 上是连续的；由于 \(\mathbf{X}_x\) 连通，\(f_x\) 必在 \(\mathbf{X}_x\) 上取常值。由归纳法立即得出，\(f(x) = f(a)\) 对一切点 \(x = (x_i)\)（其中 \(x_i = a_i\) 除有限多个指标 \(i \in I\) 外成立）都成立。但这些点 \(x\) 构成 \(\mathbf{X}\) 的一个稠密子集（\(\S 4\)，第 3 目，命题 8）。于是 \(f\) 在 \(\mathbf{X}\) 上连续，且在 \(\mathbf{X}\) 的一个稠密子集上取常值，从而在 \(\mathbf{X}\) 上取常值（\(\S 8\)，第 1 目，命题 2，推论 1）。但这与 \(f\) 的定义相矛盾。

